\section{Thiết lập thực nghiệm}
\label{sec:thiet-lap-thuc-nghiem}

\subsection{Giao thức chia dữ liệu}
\label{sec:giao-thuc-chia-du-lieu}

Dữ liệu được chia theo thời gian để mô phỏng kịch bản dự báo thực tế, tránh rò rỉ thông
tin tương lai vào quá trình chọn mô hình: \TODO{xác định cụ thể các fold rolling-origin
dùng cho validation (ví dụ huấn luyện tăng dần theo năm, validate trên năm kế tiếp) và
xác nhận tập test cuối cùng, cố định, là giai đoạn 2015--2016, không được dùng cho bất kỳ
quyết định chọn mô hình hay siêu tham số nào.}

\subsection{Baseline}
\label{sec:baseline}

\begin{itemize}
  \item \textbf{Lớp đa số:} luôn dự đoán lớp xuất hiện nhiều hơn trong tập huấn luyện.
  \item \textbf{Persistence:} dự đoán nhãn tại $t+1$ bằng nhãn quan trắc gần nhất hợp lệ
        trong cửa sổ.
  \item \textbf{Ngưỡng nhiệt độ:} phân loại Snow nếu nhiệt độ dưới một ngưỡng cố định
        (\TODO{xác định ngưỡng, ví dụ hiệu chỉnh trên tập huấn luyện}), tương ứng với
        hướng tiếp cận truyền thống ở Mục~\ref{sec:lit-phan-loai-giang-thuy}.
  \item \textbf{Logistic Regression} và \textbf{k-NN:} trên đặc trưng LiteStatTemp sau
        khi impute bằng median, làm baseline học máy đơn giản không chuyên biệt cho dữ
        liệu thiếu.
\end{itemize}

\subsection{Metrics}
\label{sec:metrics}

\TODO{Lấy định nghĩa và lý do chọn từng chỉ số từ tài liệu metrics của nhóm (ví dụ vì sao
chọn Macro F1 thay vì Accuracy do mất cân bằng lớp, và vì sao bổ sung PR-AUC để đánh giá
độ nhạy với ngưỡng quyết định). Liệt kê đầy đủ: Macro F1, Precision/Recall theo lớp,
PR-AUC, và bất kỳ chỉ số phụ nào khác (ví dụ Balanced Accuracy, Cohen's Kappa).}

\subsection{Kiểm định thống kê}
\label{sec:kiem-dinh-thong-ke}

Sự khác biệt hiệu năng giữa từng cặp mô hình trên cùng tập test được kiểm định bằng
McNemar's test \cite{mcnemar1947}, phù hợp để so sánh hai bộ phân loại trên cùng một tập
dữ liệu. Để so sánh đồng thời nhiều mô hình trên nhiều cấu hình thiếu dữ liệu, kiểm định
Friedman kèm biểu đồ critical difference (CD) \cite{demsar2006statistical} được dùng để
xếp hạng và xác định nhóm mô hình không khác biệt có ý nghĩa thống kê.

\subsection{Môi trường chạy}
\label{sec:moi-truong-chay}

\TODO{Ghi lại phần cứng (CPU/GPU, RAM), phiên bản Python và các thư viện chính (XGBoost,
sktime/aeon cho MiniRocket/MultiRocket, PyTorch cho GRU-D/BRITS/SAITS/UniTS, Optuna), và
danh sách seed cố định dùng cho khả năng tái lập kết quả.}
